%% ch09 --- Serwisy: FastAPI i pydantic dla tych, którzy znają Minimal API
%%
%% Zalążek, wygenerowany przez tools/gen_stubs.py z tools/chapters.json.
%% Poniższy opis jest kontraktem tego rozdziału: co ma uzasadnić, co ma dać
%% czytelnikowi i czego NIE ma powtarzać za tomami towarzyszącymi. Napisanie
%% rozdziału oznacza usunięcie zalążkowego bloku (tego, który zaczyna się po
%% linii \chapter) i zastąpienie go rozdziałem -- po czym ten plik nie jest
%% już generowany.

\chapter{Serwisy: FastAPI i pydantic dla tych, którzy znają Minimal API}
\label{ch:services}

\chapterstub{%
\textit{[Opis do przetłumaczenia --- patrz wydanie angielskie.]}

Argument: FastAPI is Minimal APIs with the DI container replaced by function
signatures, and pydantic is the model binder and validator you used to
configure. Routing and path operations against \code{MapGet}; \code{Depends}
against DI, with per-request lifetime by default and \code{yield}
dependencies against a disposable scope; request models, validation errors
and the 422 against \code{ProblemDetails}; response models and
serialisation, \code{model\_dump} and aliases against
\code{JsonPropertyName}; the generated OpenAPI against Swashbuckle;
middleware against ASP.NET middleware, and its order; \code{pydantic-
settings} against \code{IOptions} and configuration providers; uvicorn
against Kestrel, workers against threads, and the reverse proxy in front;
lifespan against \code{IHostedService}; background tasks; testing with
\code{TestClient} and \code{ASGITransport}. Payoff: a service skeleton the
reader can copy, and the mapping table. Measurement: experiment E5, uvicorn
workers against concurrency, throughput and p50 and p95 with a mocked
upstream, calibrated the way the LangChain book's chapter 13 recorded, free.
Five exercises. Overlap rule: the LangChain book's chapter 13 serves agents
over SSE; this chapter owns the framework mapping and hands streaming there.

\medskip
\textbf{Budżet słów:} 3000.\\
\textbf{Wydanie:} v0.2.\\
\textbf{Pomiar:} E5.
}
